\begin{center}\begin{tikzpicture}[font=\small,>=stealth,line join=round]
\draw[gray!40,rounded corners=3pt] (0,0) rectangle (14.2,7.4);
\node[text width=12.6cm,align=center] at (7.1,6.75) {为什么要相减：看输入在数轴上的距离};

\draw[->,gray] (.8,3.5)--(13.1,3.5) node[right]{$u$};
\foreach \u/\lab in {1.4/x,5.4/{x+2},11.4/{x+5}}{
\draw (\u,3.4)--(\u,3.6);\node[below=6pt] at (\u,3.4){$\lab$};}
\draw[<->,teal,thick] (1.4,4.2)--(5.4,4.2) node[midway,above]{$2$};
\draw[<->,blue,thick] (5.4,4.2)--(11.4,4.2) node[midway,above]{$5-2=3$};
\draw[<->,violet] (1.4,5.3)--(11.4,5.3) node[midway,above]{$5$};
\node[text width=11.8cm,align=center] at (7.1,1.9) {输入相距 $3$：先变号；再走 $3$：恢复。\\$5+2=7$ 把两段从 $x$ 出发的距离相加，不能表示这两点的间距。};

\node[text width=12.6cm,align=center] at (7.1,.65) {数轴按比例绘制；$f(x+5)=-f(x+2)$ 给出的一个正周期为 $6$。};
\end{tikzpicture}\end{center}
